\chapter{John Pardon}\label{ch2026johnpardon:john-pardon-fields-medal-2026}

\begin{center}\rule{0.5\linewidth}{0.5pt}\end{center}

\textbf{Prepared:} July 2026 \textbf{Subject:} John Pardon (b. 1989, USA) \textbf{Position:} Professor, Simons Center for Geometry and Physics, Stony Brook University \textbf{Doctoral Advisor:} Yakov Eliashberg (Stanford University) \textbf{Fields Medal Citation:} \emph{``For his achievements in symplectic geometry including new approaches to virtual fundamental cycles, Fukaya categories of certain manifolds and counting holomorphic curves, and for his contributions to other areas of geometry and topology, including group actions on 3-manifolds and knot theory.''}

\begin{center}\rule{0.5\linewidth}{0.5pt}\end{center}

\begin{prereqs} Differential topology (manifolds, vector bundles), symplectic geometry (symplectic manifolds, Lagrangian submanifolds), basic algebraic topology (homology, cohomology), knot theory basics. \end{prereqs}

\begin{objectives} Explain the concept of knot distortion and its significance; describe Fukaya categories and their role in symplectic geometry; understand the problem of counting holomorphic curves and virtual fundamental cycles \cite{pardon2016}. \end{objectives}

\section{Problem History}

Knot theory studies the embedding of circles in three-dimensional space. A fundamental question in metric geometry, posed by Mikhail Gromov, concerns \textbf{knot distortion}: for a knotted curve in $\mathbb{R}^3$, the distortion is the supremum over all pairs of points on the curve of the ratio of their shorter arc-length distance to their ambient Euclidean distance. A round circle has distortion $\pi/2$, while more complex knots can have higher distortion. Gromov asked whether there exist knots with arbitrarily large distortion---a question that had been open for 25 years when Pardon solved it as an undergraduate \cite{pardon2011}.

The \textbf{Hilbert--Smith conjecture} asks whether a locally compact group that acts faithfully by homeomorphisms on a manifold must be a Lie group. Equivalently, every locally compact, totally disconnected topological group acting faithfully on a manifold is discrete. This is one of the fundamental questions in geometric topology, connecting group theory, topology, and dynamics. The conjecture had been open since the 1930s.

The Hilbert--Smith conjecture is closely related to the question of whether totally disconnected groups can act ``wildly'' on manifolds. The $p$-adic integers $\mathbb{Z}_p$ provide the key test case: the conjecture is equivalent to the statement that $\mathbb{Z}_p$ cannot act faithfully by homeomorphisms on any manifold. If such an action existed, it would produce a continuous homomorphism from $\mathbb{Z}_p$ to the homeomorphism group of the manifold, which is a ``wild'' object that cannot arise from a Lie group action.

The difficulty of the conjecture lies in the fact that homeomorphisms can be much more flexible than diffeomorphisms. Smooth actions are more constrained by differential and analytic tools, while the group of homeomorphisms is much larger and has no comparable finite-dimensional Lie-group structure. Pardon's breakthrough was to show that for 3-manifolds, the extra flexibility of homeomorphisms is not enough to allow a faithful action of $\mathbb{Z}_p$.

In symplectic geometry, the \textbf{Gromov--Witten invariants} are defined by counting pseudo-holomorphic curves in a symplectic manifold $(M,\omega)$ satisfying boundary or incidence conditions. A pseudo-holomorphic curve is a map $u:(\Sigma,j)\to(M,\omega)$ satisfying the Cauchy--Riemann equation $\bar{\partial}_J u = 0$. However, the moduli space of such curves is often singular or of the wrong dimension, making a naive count impossible.

The standard approach is to construct a \textbf{virtual fundamental cycle}: a homology class in the moduli space that plays the role of the fundamental class if the space were smooth and of the expected dimension. Previous constructions (by Li--Tian, Behrend--Fantechi, Fukaya--Ono, and others) worked in specific settings but had limitations.

\textbf{Fukaya categories} are $A_\infty$-categories associated to a symplectic manifold $(M,\omega)$, whose objects are Lagrangian submanifolds $L \subset M$ and whose morphisms $\mathrm{hom}(L_0,L_1)$ are generated by intersection points, with composition defined by counting pseudo-holomorphic disks. They play a central role in homological mirror symmetry (Kontsevich, 1994), connecting symplectic geometry to algebraic geometry through the mirror symmetry correspondence.

However, foundational issues---particularly in the construction of virtual fundamental cycles for moduli spaces of holomorphic curves in Liouville manifolds (exact symplectic manifolds with contact-type boundary)---had remained unresolved. These manifolds arise naturally as complements of hypersurfaces in projective space and as Stein manifolds.

\section{How Was It Solved}

As an undergraduate at Princeton, Pardon solved Gromov's distortion problem by proving that for any constant $C$, there exist knots whose distortion exceeds $C$. The proof constructs explicit families of knots with controlled geometry, showing that the ratio of arc-length to ambient distance can be made arbitrarily large using techniques from geometric measure theory and the analysis of embeddings.

Pardon proved the Hilbert--Smith conjecture for ``rough'' (non-differentiable) group actions on 3-manifolds: if a totally disconnected, locally compact group acts faithfully by homeomorphisms on a 3-manifold, then the group must be discrete. This was the first case of the Hilbert--Smith conjecture proved for actions on manifolds of any dimension. The proof uses contact-local isotopy techniques and the theory of group actions on low-dimensional manifolds.

He also proved Edmonds' smoothing conjecture for rough group actions on 3-manifolds, showing that such actions can be approximated by smooth actions. This complements the Hilbert--Smith result and provides insight into the relationship between topology and dynamics.

For symplectic geometry, Pardon developed an algebraic framework for constructing virtual fundamental cycles on moduli spaces of pseudo-holomorphic curves. His central tool is an \emph{implicit atlas}: a compatible system of local finite-dimensional reductions that can be assembled into a global virtual cycle. The construction uses topological gluing results and does not rely on perturbing the defining equations \cite{pardon2016}. Pardon also applied this framework to define Gromov--Witten invariants and Hamiltonian Floer homology over $\mathbb{Q}$.

This framework addresses the failure of transversality in holomorphic-curve moduli spaces, where singularities and excess dimension prevent a direct count. Pardon’s work contributes to the foundations for counting holomorphic curves and studying Fukaya categories; specific applications should be tied to the relevant results rather than attributed to a derived-stack construction.

\section{Impact of This Discovery}

Pardon's work on knot distortion settled a fundamental question in metric geometry, showing that the geometry of knotted curves in three-dimensional space can be far more complex than naive intuition suggests. This result connects knot theory to bi-Lipschitz embeddings and coarse geometry.

The Hilbert--Smith conjecture for 3-manifolds is a major advance in geometric topology. While the full conjecture in higher dimensions remains open, Pardon's result demonstrates that the conjecture is true in the lowest interesting dimension and provides techniques that may be applicable more broadly.

The new construction of virtual fundamental cycles has transformed the study of Gromov--Witten invariants and Fukaya categories. It provides a rigorous foundation for enumerative geometry in symplectic settings that were previously inaccessible, with applications to homological mirror symmetry and the study of Lagrangian submanifolds.

Pardon's contributions to the foundations of Fukaya categories have resolved long-standing foundational issues that had hindered progress in symplectic topology. His work has enabled new results on wrapped Floer theory, Lagrangian embeddings, and the symplectic topology of Liouville manifolds.

His methods have been adopted by researchers worldwide and have opened new directions of inquiry in symplectic geometry, particularly in the study of moduli problems with non-trivial obstructions.

The \textbf{Gromov--Witten invariants} are defined by counting pseudo-holomorphic curves in a symplectic manifold $(M, \omega)$ satisfying boundary or incidence conditions. For a compact symplectic manifold of dimension $2n$, the Gromov--Witten invariant of class $A \in H_2(M)$ and genus $g$ is defined by
\[
\langle \gamma_1, \ldots, \gamma_k \rangle_{g,A} = \int_{[\overline{\mathcal{M}}_{g,k}(M,A)]^{vir}} \mathrm{ev}_1^*(\gamma_1) \wedge \cdots \wedge \mathrm{ev}_k^*(\gamma_k),
\]
where $\overline{\mathcal{M}}_{g,k}(M,A)$ is the moduli space of stable maps of genus $g$ with $k$ marked points in class $A$, and $[\overline{\mathcal{M}}_{g,k}(M,A)]^{vir}$ is the virtual fundamental cycle. The virtual fundamental cycle is needed because the moduli space is typically singular and of the wrong dimension, making a naive count impossible.

Pardon's 2016 paper develops the implicit-atlas method for virtual fundamental cycles and illustrates it by defining Gromov--Witten invariants and Hamiltonian Floer homology over $\mathbb{Q}$ \cite{pardon2016}. This construction supplies a framework for obtaining invariants from moduli spaces that are not cut out transversally. It is distinct from constructing a derived stack with a perfect obstruction theory.

\section{Modern Views and Improvements}

Major open questions include the full Hilbert--Smith conjecture in higher dimensions, the construction of Fukaya categories for more general symplectic manifolds, and the development of new invariants from the count of holomorphic curves.

Pardon's virtual fundamental cycle construction has been extended and refined by other researchers, leading to new applications in enumerative geometry and the study of Donaldson--Thomas type invariants in symplectic settings.

Virtual-cycle constructions remain an active area of symplectic geometry. Pardon's implicit-atlas framework is one approach to defining invariants from moduli spaces that fail to be transverse.

Future directions include the development of Floer-theoretic invariants for more exotic symplectic manifolds, the application of Pardon's techniques to contact topology and Legendrian knots, and the exploration of connections between symplectic geometry and quantum field theory.

The full Hilbert--Smith conjecture in higher dimensions remains one of the most important open problems in geometric topology. While Pardon's result for 3-manifolds is the strongest known case, the conjecture in dimensions 4 and above requires fundamentally new techniques. The difficulty is that in higher dimensions, homeomorphisms are even more flexible than in dimension 3, making it harder to rule out wild group actions.

The study of Fukaya categories for Liouville manifolds is a broader research program. Pardon’s work on virtual fundamental cycles contributes to its foundations, while particular computations and mirror-symmetry results require their own sources.

Pardon's results on group actions on 3-manifolds and his work in symplectic geometry are distinct contributions; the chapter does not infer a direct method transfer between them.

Later work on derived algebraic geometry provides other approaches to moduli problems and virtual classes. These should be distinguished from Pardon's implicit-atlas construction.

The study of knot invariants, motivated by Pardon's work on knot distortion, continues to be an active area. The distortion of a knot is a metric invariant that measures the ratio of arc-length distance to ambient Euclidean distance, and it provides information about the geometry of the knot. The study of knot distortion and related invariants has connections to the theory of bi-Lipschitz embeddings, to geometric measure theory, and to the study of fractal geometry in three dimensions.

Donaldson--Thomas theory and contact topology are related neighboring subjects, but claims that Pardon’s constructions directly established results in these areas require separate references and are not made here.

An exact Lagrangian is a Lagrangian submanifold on which the symplectic primitive restricts to an exact one-form. Fukaya categories provide a framework for studying Lagrangians and their intersections. Specific classification results are not attributed to Pardon here without direct references.

The \textbf{wrapped Fukaya category} extends Fukaya-category constructions to suitable noncompact Lagrangians in exact symplectic manifolds. Its definitions and applications are a broader program; this chapter does not attribute the entire theory to Pardon.

Connections between symplectic geometry and quantum field theory are an independent research direction and are not a direct consequence of the results discussed here.

Categorical invariants remain an active direction in symplectic geometry. Their development is broader than the specific virtual-cycle construction discussed in this chapter.

\begin{center}\rule{0.5\linewidth}{0.5pt}\end{center}

\section*{Summary}
\addcontentsline{toc}{section}{Summary}

Pardon solved Gromov's knot distortion problem, proved the Hilbert--Smith conjecture for 3-manifolds, and developed new foundations for virtual fundamental cycles and Fukaya categories in symplectic geometry.
